\section{问题建模}

\subsection{任务场景与海洋再分析数据}
本文研究自主水下航行器在强流海区携带陈旧粗分辨率海流图执行绕障航行的二元路径决策问题，全部环境设置与参数汇总于表\ref{tab:settings}。研究海区选取台湾以东至东海的黑潮强流区（23°--31°N、121°--129°E），海流数据取自HYCOM全球海洋再分析产品（GLBy0.08，实验93.0\cite{chassignet2007hycom}）10米水深处的水平流速场，空间分辨率为1/12°，时间跨度自2019年1月1日至2024年8月29日每3天取一张12:00时快照，共689张；其中2019--2021年共365张作为开发年份用于标定误差模型与门控阈值，2022--2024年共324张作为独立测试年份。在每张海流快照上随机生成2个航行任务，要求起点与终点相距100--140千米且离岸不少于20千米，两点连线中点处存在一个半径15千米的圆形禁航区，航行器必须选择从禁航区的左侧或右侧绕过；为聚焦强流影响，只保留粗图上任务区平均流速不低于0.3米/秒的任务，最终得到开发年份710个任务（分布于357天）与测试年份623个任务（分布于313天）。航行器在下水时仅携带3天前生成的1°粗分辨率海流图（由3天前的1/12°场在1°网格内平均得到，简称粗图），而当天的1/12°再分析海流场则作为检验决策好坏的真实海流。

\begin{table}[htbp]
\centering
\caption{粗分辨率海流图下自主水下航行器航路决策问题的实验设置与参数。}
\label{tab:settings}
\small
\begin{tabular}{@{}p{0.16\linewidth}p{0.80\linewidth}@{}}
\toprule
项目 & 设置 \\
\midrule
海流数据 & HYCOM 全球海洋再分析（GLBy0.08，实验 93.0），水平分辨率 1/12°，取 10 米深度；23°--31°N、121°--129°E（台湾以东至东海，含黑潮）；2019-01-01 至 2024-08-29 每 3 天 12:00 时一张，共 689 张 \\
任务 & 每张快照抽 2 个：起终点相距 100--140 千米、离岸不少于 20 千米，两点中间设半径 15 千米的禁航区，可从左侧或右侧绕过；只取粗图上任务区平均流速不低于 0.3 米/秒的强流任务；开发 710 个（2019--2021 年），测试 623 个（2022--2024 年） \\
粗图与真实海流 & 航行器带 3 天前的 1° 粗分辨率海流图（下称粗图，格内平均）；真实海流为当天 1/12° 再分析海流场 \\
航行器 & 对水速度 1.0 米/秒；选定一侧后，按真实海流在该侧以最快方式航行 \\
航线求解 & 网格最短路：粗网格求解 8 千米网格、8 个方向；中网格求解 4 千米、16 个方向；细网格求解 2 千米、32 个方向（计算量之比 1 : 8 : 64）；顶流过强的边不可通行 \\
合理流场集合 & 粗图加上由流函数生成的无散误差模态的组合，各点误差大小与相关长度（76 千米）由开发年份标定，组合系数服从标准正态；半径 $r=1.5$（开发年份中包住真实左右航时差的任务占 93\%） \\
合理反例检验 & 在半径 $r$ 的球内对左右航时差做投影梯度上升（梯度沿最快航线累加）；两个方向、半径 1 与 2 各搜一次，其他半径线性插值，每个任务约 71 次中网格求解 \\
观测方式 & 原地停留测流（出发前 2 小时测 12 次）；沿途测流 30 千米或沿途测流 60 千米后可换边；绕行定点测流（去前方 40 千米内反例与粗图差别最大处测量）；上浮取当天 1° 图或上浮取当天 0.5° 图（每次 1 小时）；测流噪声 0.05 米/秒 \\
阈值判据 & 最大后悔值超过 8 小时（开发年份选定）时判定粗图不可信，触发超阈值时保守改选或超阈值时上浮取图，否则按粗图选边 \\
评价指标 & 额外航时 = 实际航时 $-$ 真实更快一侧的航时（小时）；配对差值的 95\% 置信区间按天整组自助重抽样 2000 次 \\
\bottomrule
\end{tabular}

\end{table}

\subsection{航线求解与计算预算}
给定任一水平海流场 $\boldsymbol{u}(\boldsymbol{x})$ 与航行器恒定的对水巡航速度 $V = 1.0$ 米/秒，航行器在选定绕行一侧 $s \in \{\text{L}, \text{R}\}$（左侧或右侧）后的最快航线通过有向图上的最短路算法求取\cite{dijkstra1959note}。将任务海域离散为规则空间网格，对连接相邻网格节点的有向边 $e$（长度为 $\ell_e$、单位方向向量为 $\boldsymbol{d}_e$），以边两端节点海流的平均作为该边的海流，记其沿边分量为 $u_{\parallel} = \boldsymbol{u} \cdot \boldsymbol{d}_e$、垂直边分量为 $u_{\perp} = \|\boldsymbol{u} - u_{\parallel}\boldsymbol{d}_e\|$。为保持航迹严格沿边方向推进，航行器必须以部分对水速度抵消横流，由此得到沿该边的有效对地速度 $v_{\text{g}}(e)$ 与通行时间 $t(e)$ 分别为：
\begin{equation}
v_{\text{g}}(e) = u_{\parallel} + \sqrt{V^2 - u_{\perp}^2}, \qquad t(e) = \frac{\ell_e}{v_{\text{g}}(e)}.
\label{eq:edge_time}
\end{equation}
当横流过大导致 $|u_{\perp}| \ge V$ 或顶流过强导致 $v_{\text{g}}(e) \le 0$ 时，该有向边不可通行，其通行时间记为无穷大。通过在连接起点与终点且严格位于禁航区左侧或右侧的子图上求解最短路径，即可得到流场 $\boldsymbol{u}$ 下左、右两侧的最优航时 $T_{\text{L}}(\boldsymbol{u})$ 与 $T_{\text{R}}(\boldsymbol{u})$，并定义两侧时间差为：
\begin{equation}
\Delta T(\boldsymbol{u}) = T_{\text{L}}(\boldsymbol{u}) - T_{\text{R}}(\boldsymbol{u}),
\label{eq:delta_t}
\end{equation}
其中 $\Delta T(\boldsymbol{u}) < 0$ 表示左侧更快，$\Delta T(\boldsymbol{u}) > 0$ 表示右侧更快。为刻画机载计算资源的投入，本文设置三档航线求解精度：粗算采用8千米网格与8个航向（相对计算量为1），标准采用4千米网格与16个航向（相对计算量为8），精算采用2千米网格与32个航向（相对计算量为64，按图中有向边总数计量）。选定绕行一侧 $s$ 后，假设航行器在实际下水航行过程中能够利用走航感知在该侧拓扑通道内按真实海流 $\boldsymbol{u}_{\text{true}}$ 以最优方式顺应流场（在测试年份中更快一侧的平均航时为39.2小时），因此本文聚焦于出发前或航行初段“选左侧还是选右侧”这一不可逆的二元拓扑决策。

\subsection{合理海流集合的参数化建模}
由于3天前的1°粗图 $\boldsymbol{u}_0(\boldsymbol{x})$ 同时缺失了中尺度切变与3天内的流轴演变，真实海流必然偏离 $\boldsymbol{u}_0$，但其偏离幅度与空间结构受区域海洋动力学约束而非任意噪声。我们用流函数生成无散的误差模态，并按多年数据标定的误差大小逐点缩放，构造叠加在粗图之上的合理海流 $\boldsymbol{u}(\boldsymbol{x}; \boldsymbol{z})$：
\begin{equation}
\boldsymbol{u}(\boldsymbol{x}; \boldsymbol{z}) = \boldsymbol{u}_0(\boldsymbol{x}) + \sum_{k=1}^{K} z_k \boldsymbol{m}_k(\boldsymbol{x}), \qquad \boldsymbol{m}_k(\boldsymbol{x}) = \mathbf{S}(\boldsymbol{x})\, \nabla^{\perp}\phi_k(\boldsymbol{x}),
\label{eq:plausible_field}
\end{equation}
其中 $\nabla^{\perp} = (-\partial_y, \partial_x)$，$\{\phi_k\}_{k=1}^{K}$ 为相关长度 $L_{\text{corr}}$ 的高斯相关流函数的前 $K$ 个主模态（保留95\%方差），因而 $\nabla^{\perp}\phi_k$ 处处无散；$\mathbf{S}(\boldsymbol{x})$ 为对角缩放矩阵，使两个流速分量在各点的方差等于开发数据标定的误差方差（缩放后近似无散）；$\boldsymbol{z} = (z_1, \dots, z_K)^{\top}$ 为先验服从标准正态分布的组合系数。给定半径 $r > 0$，合理集合定义为
\begin{equation}
\mathcal{U}(r) = \left\{ \boldsymbol{u}(\boldsymbol{x}; \boldsymbol{z}) \;\middle|\; \|\boldsymbol{z}\|_2 \le r \right\}.
\label{eq:plausible_set}
\end{equation}
半径 $r$ 由开发年份标定（第5.3节），使 $\mathcal{U}(r)$ 恰好大到能包住大多数真实海洋的两侧时间差。

\subsection{可选动作、观测代价与评价指标}
在任务起点处，航行器面对粗图 $\boldsymbol{u}_0$ 可在五类动作中做出选择：（1）直接照地图走：在粗图上用粗算、标准或精算选出账面更快的一侧立即出发；（2）改变决策准则而不取新信息：在合理集合上按期望时间选平均航时更短的一侧，或采取保守准则选最坏情况下后悔最小的一侧；（3）就地或顺路原位测流：包括出发前原地停留2小时连续测流12次、沿所选一侧航行30千米或60千米并连续测流后再决定是否调头换边，以及专门绕行到前方40千米内反例与地图流速差别最大处去定点测量（原位测流仪器噪声为0.05米/秒，执行数据同化更新流场时考虑代表性误差取0.15米/秒\cite{anderson2001eakf}）；（4）上浮取新图：上浮接收岸基下发的当天1°图或当天0.5°图，每次上浮通信与定位耗时计1小时（并在灵敏度分析中考察0至4小时的代价变化），随后在更新后的地图上选边出发；（5）本文提出的反例检验门控：先进行按需精算与反例检验，仅在最坏后悔超标时触发保守选边或上浮取图。设某决策动作最终耗费的总航行与观测时间为 $T_{\text{total}}$，定义该次任务的航时“损失”为实际总耗时减去真实海流下更快一侧的最优航时：
\begin{equation}
\mathcal{L} = T_{\text{total}} - \min\left\{ T_{\text{L}}(\boldsymbol{u}_{\text{true}}),\, T_{\text{R}}(\boldsymbol{u}_{\text{true}}) \right\}.
\label{eq:loss}
\end{equation}
损失 $\mathcal{L} \ge 0$ 越低越好；当不发生额外驻留、绕行或上浮耗时且选对更快一侧时，损失为0小时。

\subsection{问题剖析：为何常规计算与盲目观测难以奏效}
上述决策问题的核心矛盾在于：选对左右绕行一侧高度依赖禁航区两侧数十千米尺度上的真实海流切变信息，而航行器手中的3天前1°粗图恰恰抹平了这一信息，机载计算只能不断锐化粗图自身的几何推论，却无法凭空创造缺失的海洋物理信息。这一根本矛盾衍生出两个决定方法设计的结构性性质：

第一，粗图上的两侧时间差模糊性具有两个截然不同的来源——求解器网格离散引起的“数值歧义”与海洋背景场失真引起的“环境歧义”，二者要求的应对动作完全不同。当 $|\Delta T(\boldsymbol{u}_0)|$ 很小时，如果仅是因为8千米粗算网格过粗导致路径折线失真，换成2千米精算网格即可看清粗图本来的倾向；然而，如果即使在精算下地图选出了某侧 $s_0$，在合理集合 $\mathcal{U}(r)$ 内依然存在另一个物理相容的流场 $\boldsymbol{u}^*$ 使另一侧明显更快，那么这种模糊性就是环境本身的多解性，继续加密网格到1千米或0.5千米也丝毫无助于判断真实海洋究竟接近 $\boldsymbol{u}_0$ 还是 $\boldsymbol{u}^*$。

第二，是否值得付出宝贵的航时去执行观测，并不取决于是否存在合理反例，而取决于该反例一旦成真会让当前决策多损失多少小时（即最坏后悔的大小）以及观测手段能否触及决定胜负的海区。在强流区使用3天前的1°粗图时，由于背景误差相关长度达数十千米，几乎任何任务都能在合理集合内找到某个让另一侧略快一点的流场；若一见到反例就去观测或放弃地图，系统将陷入过度反应。因此，算法必须直接求出合理集合内所选一侧相对于另一侧的最大可能航时劣势，并将这一最坏后悔与观测的时间代价直接比较，这自然引出了第5节的合理反例检验与门控方法。
